\documentclass[12pt,a4paper]{article}
\usepackage[utf8]{inputenc}
\usepackage[T5]{fontenc}
\usepackage[vietnamese]{babel}
\usepackage{geometry}
\geometry{margin=2.2cm}
\usepackage{amsmath,amssymb,amsthm}
\usepackage{booktabs,longtable,tabularx,multirow}
\usepackage{graphicx}
\usepackage{hyperref}
\usepackage{listings}
\usepackage{xcolor}
\usepackage{enumitem}
\usepackage{float}
\usepackage{titlesec}
\titleformat{\section}{\large\bfseries}{\thesection}{0.5em}{}
\titleformat{\subsection}{\normalsize\bfseries}{\thesubsection}{0.5em}{}

\lstset{basicstyle=\ttfamily\small,breaklines=true,frame=single,backgroundcolor=\color{gray!8}}

\title{BÁO CÁO ĐẦY ĐỦ: Planning Optimization\\
\large Từ nền tảng lý thuyết đến triển khai thực nghiệm đa siêu tham số}
\author{Tổng hợp và triển khai trên mã nguồn dự án}
\date{\today}

\begin{document}
\maketitle
\tableofcontents
\newpage

\section{Giới thiệu}
Bài toán trong dự án là lập kế hoạch thu hoạch nông sản với ba nhóm ràng buộc chính: (i) mỗi thửa ruộng chỉ được thu hoạch trong cửa sổ thời gian $[s_i,e_i]$, (ii) tổng sản lượng xử lý mỗi ngày không vượt quá công suất nhà máy $M$, và (iii) nếu một ngày hoạt động thì tổng tải phải đạt ngưỡng tối thiểu $m$. Mục tiêu là tối đa tổng sản lượng được thu hoạch.

Về mặt tối ưu hóa tổ hợp, đây là một bài toán khó do vừa có biến nhị phân/chọn không chọn, vừa có biến gán ngày, vừa có ràng buộc logic ``hoạt động hoặc không hoạt động'' theo ngày. Vì vậy dự án triển khai nhiều hướng: Greedy, Local Search, Dynamic Programming (2 biến thể), MIP và CP-SAT, sau đó benchmark toàn diện theo nhiều bộ siêu tham số.

\subsection{Mục tiêu báo cáo}
Báo cáo này mở rộng từ bản tóm tắt \texttt{planning\_optimization\_report.md} thành bản đầy đủ, với các mục tiêu:
\begin{itemize}[leftmargin=1.4cm]
  \item Giải thích lại chi tiết các thuật toán ở mức người mới bắt đầu vẫn hiểu.
  \item Trình bày tư duy mô hình hóa từng nhóm phương pháp.
  \item Chạy thực nghiệm với nhiều siêu tham số và phân tích kết quả.
  \item Đưa ra khuyến nghị chọn phương pháp theo bối cảnh vận hành thực tế.
\end{itemize}

\section{Mô tả bài toán và ký hiệu}
Dữ liệu đầu vào gồm $N$ thửa ruộng, mỗi ruộng $i$ có:
\begin{itemize}
  \item sản lượng $d_i$,
  \item ngày bắt đầu thu hoạch sớm nhất $s_i$,
  \item ngày kết thúc thu hoạch muộn nhất $e_i$.
\end{itemize}
Gọi $x_i \in \{-1,1,2,\dots,T\}$, trong đó $x_i=-1$ nghĩa là không thu hoạch ruộng $i$, ngược lại $x_i=t$ nghĩa là thu hoạch ở ngày $t$.

Ràng buộc theo ngày $t$:
\begin{equation}
L_t = \sum_{i: x_i=t} d_i,\qquad L_t \le M,
\end{equation}
\begin{equation}
L_t=0\;\text{hoặc}\;L_t\ge m.
\end{equation}
Ràng buộc cửa sổ:
\begin{equation}
x_i=t>0 \Rightarrow s_i \le t \le e_i.
\end{equation}
Mục tiêu:
\begin{equation}
\max \sum_{i: x_i>0} d_i.
\end{equation}

\section{Nền tảng lý thuyết (mức beginner)}
\subsection{Linear Programming (LP)}
LP là tối ưu một hàm tuyến tính với ràng buộc tuyến tính. Điểm mấu chốt cho người mới: miền nghiệm LP là một đa diện lồi, và nghiệm tối ưu (nếu hữu hạn) nằm ở đỉnh. Vì vậy solver kiểu simplex đi ``nhảy từ đỉnh này qua đỉnh khác'' để cải thiện mục tiêu.

Trong bài toán này, LP thường được dùng như \textbf{relaxation} (thư giãn) cho mô hình nguyên: cho phép biến nhận giá trị phân số để lấy cận trên và định hướng nhánh-cận.

\subsection{Mixed Integer Programming (MIP)}
MIP là LP + yêu cầu một số biến phải nguyên/nhị phân. Trực giác: biến nhị phân giúp mô hình hóa quyết định có/không; biến nguyên giúp mô hình hóa số lượng, trạng thái theo ngày. Nhưng đổi lại, độ khó tăng mạnh (thường NP-hard).

Solver MIP hiện đại dùng ba trụ cột:
\begin{enumerate}
  \item \textbf{LP relaxation} để lấy cận,
  \item \textbf{Branch-and-Bound} để chia không gian nghiệm,
  \item \textbf{Cutting planes/heuristics} để tăng tốc.
\end{enumerate}

\subsection{Constraint Programming và CP-SAT}
CP tập trung vào ràng buộc logic và propagation miền giá trị. CP-SAT trong OR-Tools là lai giữa SAT + CP + linear constraints nên rất mạnh cho bài toán có nhiều cấu trúc tổ hợp.

Đối với người mới: hãy nghĩ CP-SAT như một ``máy suy luận ràng buộc'' rất giỏi xử lý điều kiện if-then, ExactlyOne, implication, reification và các ràng buộc nguyên.

\subsection{Heuristic và Metaheuristic}
Khi dữ liệu lớn, heuristic thường cho nghiệm tốt rất nhanh:
\begin{itemize}
  \item \textbf{Greedy}: quyết định cục bộ nhanh.
  \item \textbf{Local Search}: bắt đầu từ nghiệm ban đầu, lặp cải tiến theo các phép biến đổi lân cận.
  \item \textbf{Hybrid Greedy + Repair + Fill}: kết hợp xây dựng nghiệm và sửa nghiệm vi phạm.
\end{itemize}

\section{Mô hình hóa bài toán theo từng trường phái}
\subsection{MIP formulation cơ bản}
Khai báo biến nhị phân:
\begin{align}
 y_{i,t} &= 1 \Leftrightarrow \text{ruộng }i\text{ thu hoạch ngày }t,\\
 z_t &= 1 \Leftrightarrow \text{ngày }t\text{ nhà máy hoạt động}.
\end{align}

Ràng buộc:
\begin{align}
&\sum_{t=s_i}^{e_i} y_{i,t} \le 1, \quad \forall i,\\
&L_t = \sum_i d_i y_{i,t},\\
&L_t \le M z_t,\\
&L_t \ge m z_t,\\
&y_{i,t}\in\{0,1\}, z_t\in\{0,1\}.
\end{align}
Mục tiêu:
\begin{equation}
\max \sum_i\sum_{t=s_i}^{e_i} d_i y_{i,t}.
\end{equation}

\paragraph{Giải thích beginner:} $z_t$ đóng vai trò ``công tắc''. Nếu $z_t=0$ thì tải ngày đó buộc về 0; nếu $z_t=1$ thì tải phải nằm trong $[m,M]$.

\subsection{CP-SAT formulation}
Cùng tư duy biến như MIP nhưng biểu diễn bằng IntVar/BoolVar, thêm các ràng buộc implication. CP-SAT thường nhanh khi có nhiều ràng buộc logic phức tạp.

\subsection{Greedy + Repair + Fill}
\textbf{Greedy phase}: gán các ruộng theo thứ tự ưu tiên (cửa sổ hẹp trước, deadline sớm trước, sản lượng lớn trước).\\
\textbf{Repair phase}: nếu có ngày tải nằm trong $(0,m)$ thì di chuyển/bỏ bớt ruộng để loại vi phạm.\\
\textbf{Fill phase}: cố gắng chèn ruộng chưa chọn vào ngày đang hoạt động hoặc mở ngày mới khi có thể đạt ngưỡng $m$.

\subsection{Local Search (các move)}
Các move chính trong biến thể local search:
\begin{itemize}
  \item \textbf{insert}: chèn ruộng chưa chọn vào ngày khả thi,
  \item \textbf{relocate}: dời ruộng từ ngày A sang B,
  \item \textbf{swap}: đổi ngày của hai ruộng,
  \item \textbf{replace}: thay ruộng nhỏ bằng ruộng lớn chưa chọn,
  \item \textbf{group}: dời nhóm ruộng nhỏ giữa các ngày.
\end{itemize}

\section{Thuật toán chi tiết (kèm pseudo-code)}
\subsection{Greedy chi tiết}
\begin{lstlisting}
Input: danh sách ruộng (d_i, s_i, e_i), m, M
Sắp thứ tự ruộng theo (độ rộng cửa sổ tăng, deadline tăng, sản lượng giảm)
for mỗi ruộng i:
    tìm ngày t khả thi tốt nhất (ưu tiên ngày đã có tải cao nhất nhưng chưa vượt M)
    nếu có t thì gán i vào t
Lặp REPAIR_FILL_ROUNDS lần:
    repair()  # xử lý ngày vi phạm 0 < load < m
    fill()    # chèn thêm ruộng chưa chọn để tăng tổng sản lượng
Xuất nghiệm
\end{lstlisting}

\subsection{Local Search chi tiết}
Ý tưởng quan trọng cho người mới: local search không ``sinh nghiệm từ đầu'' mà cải thiện quanh nghiệm hiện tại. Chất lượng phụ thuộc mạnh vào:
\begin{itemize}
  \item nghiệm khởi tạo,
  \item tập move,
  \item thứ tự duyệt ứng viên,
  \item số vòng lặp tối đa.
\end{itemize}

\subsection{DP truyền thống theo ngày}
Xem mỗi ngày như một bài knapsack cục bộ với danh sách ruộng có thể gán ngày đó. Hạn chế: quyết định cục bộ ngày trước có thể làm mất cơ hội ngày sau. Vì vậy thêm tham số \texttt{PASSES} để quét nhiều lượt.

\subsection{DP local-knapsack hybrid}
Bắt đầu từ greedy, sau đó ở từng ngày chạy cải tiến kiểu knapsack trong tập ứng viên giới hạn \texttt{DP\_CANDIDATE\_LIMIT}. Đây là cách cân bằng tốc độ/chất lượng.

\subsection{MIP và CP trong dự án}
Hai mô hình exact đều dùng chiến lược \textbf{candidate days}: thay vì cho phép mọi ngày trong $[s_i,e_i]$, chỉ giữ tập ngày đại diện (ví dụ $s,e,\lfloor(s+e)/2\rfloor$ và một số điểm nội suy), giúp giảm số biến/ràng buộc và tăng tốc đáng kể.

\section{Thiết kế thực nghiệm}
\subsection{Môi trường và bộ test}
Thực nghiệm chạy bằng script \texttt{project/eval.py} trên 5 test chuẩn \texttt{project/tests/test1..test5}. Mỗi nghiệm hợp lệ và đạt giá trị mục tiêu không thấp hơn nghiệm tham chiếu được chấm 100 điểm/test.

\subsection{Lưới siêu tham số đã chạy}
Tổng cộng 42 cấu hình thực nghiệm:
\begin{itemize}
  \item Greedy: 6 cấu hình (2 luật sắp xếp $\times$ 3 mức rounds).
  \item Local-search original: 3 cấu hình rounds.
  \item Local-search variants: 12 cấu hình move/rounds/top-k.
  \item DP local: 3 mức candidate limit.
  \item DP traditional: 6 cấu hình (candidate limit $\times$ passes).
  \item MIP: 6 cấu hình (time limit + divisor).
  \item CP: 6 cấu hình (time limit + divisor).
\end{itemize}

\section{Kết quả thực nghiệm tổng hợp}
Tất cả 42 cấu hình đều đạt \textbf{500/500 điểm} trên 5 test, nghĩa là chất lượng nghiệm (theo thang chấm của benchmark) đều đạt yêu cầu. Khác biệt chính nằm ở thời gian chạy.

\subsection{Bảng thời gian theo nhóm}
\begin{longtable}{llrr}
\toprule
Nhóm & Cấu hình tiêu biểu & Điểm & Thời gian (s)\\
\midrule
\endhead
Greedy & window\_deadline\_amount-r1 & 500 & 0.253 \\
Greedy & deadline\_window-r4 & 500 & 0.294 \\
Local search & original-r5 & 500 & 0.281 \\
Local variants & all-moves (nhẹ) & 500 & 0.433 \\
Local variants & all-moves (nặng) & 500 & 0.540 \\
DP local & c160 & 500 & 0.364 \\
DP traditional & c320-p2 & 500 & 2.773 \\
MIP & time-8000 & 500 & 76.476 \\
CP-SAT & time-8.0 & 500 & 22.560 \\
\bottomrule
\end{longtable}

\subsection{Phân tích theo siêu tham số}
\paragraph{Greedy rounds:} tăng \texttt{REPAIR\_FILL\_ROUNDS} từ 1 lên 4 làm tăng nhẹ runtime, nhưng điểm không đổi.

\paragraph{Local search rounds và top-k:} cấu hình nặng (nhiều rounds/top-k lớn/all moves) tốn thời gian hơn rõ rệt, nhưng benchmark hiện tại chưa đủ khó để thể hiện chênh lệch chất lượng.

\paragraph{DP candidate limit:} tăng giới hạn ứng viên không cải thiện điểm trong bộ test hiện có, nhưng làm tăng thời gian vừa phải.

\paragraph{MIP/CP time limit:} do test tương đối ``dễ'' với mô hình candidate days, giảm time limit vẫn đạt điểm tối đa; runtime tổng phần lớn bị chi phối bởi test lớn.

\section{Giải thích dễ hiểu: vì sao mọi cấu hình đều 500/500?}
Có ba nguyên nhân:
\begin{enumerate}
  \item Bộ test tham chiếu dùng ngưỡng ``không thấp hơn reference'', không yêu cầu vượt trội tuyệt đối.
  \item Nhiều thuật toán đã tích hợp repair/fill tốt, nên dễ đạt nghiệm khả thi chất lượng cao.
  \item Candidate days giảm mạnh độ khó mô hình exact nhưng vẫn giữ đủ phương án tốt.
\end{enumerate}

\section{Khuyến nghị chọn thuật toán trong thực tế}
\begin{itemize}
  \item Cần rất nhanh (online/real-time): \textbf{Greedy hoặc Local Search nhẹ}.
  \item Cân bằng chất lượng/tốc độ: \textbf{DP local-knapsack hybrid}.
  \item Cần mô hình hóa giàu ràng buộc và nghiệm chắc chắn: \textbf{CP-SAT}.
  \item Cần diễn giải theo chuẩn OR cổ điển, khai thác dual bound rõ: \textbf{MIP}.
\end{itemize}

\section{Hướng mở rộng để tạo khác biệt chất lượng rõ hơn}
\begin{enumerate}
  \item Tạo test khó hơn: cửa sổ hẹp xen kẽ, nhiều xung đột dung lượng, phân phối $d_i$ lệch mạnh.
  \item Đổi metric: ngoài score 0/100, so thêm objective tuyệt đối và optimality gap.
  \item Với local search: thêm cơ chế chấp nhận tạm thời nghiệm xấu (simulated annealing/tabu).
  \item Với MIP/CP: warm start từ greedy để giảm thời gian hội tụ.
\end{enumerate}

\section{Kết luận}
Báo cáo đã trình bày đầy đủ từ lý thuyết đến triển khai, giải thích chi tiết từng thuật toán theo hướng beginner-friendly, và chạy benchmark đa siêu tham số trên toàn bộ pipeline. Kết quả cho thấy chất lượng nghiệm rất ổn định trên bộ test hiện tại; sự khác biệt chủ yếu nằm ở thời gian chạy: heuristic/DP nhanh hơn rất nhiều so với MIP, trong khi CP-SAT nằm giữa.

\appendix
\section{Một số cấu hình siêu tham số đã dùng}
\begin{itemize}
  \item \textbf{Greedy}: \texttt{GREEDY\_REPAIR\_FILL\_ROUNDS \{1,2,4\}}, \texttt{GREEDY\_ORDER\_MODE \{window\_deadline\_amount, deadline\_window\}}.
  \item \textbf{Local variants}: \texttt{LS\_MOVES}, \texttt{LS\_MAX\_ROUNDS}, \texttt{LS\_TOP\_ASSIGNED}, \texttt{LS\_TOP\_UNASSIGNED}, \texttt{LS\_GROUP\_DAYS}, \texttt{LS\_GROUP\_SIZE}.
  \item \textbf{DP}: \texttt{DP\_CANDIDATE\_LIMIT \{80,160,320\}}, \texttt{DP\_TRAD\_PASSES \{1,2\}}.
  \item \textbf{MIP}: \texttt{MIP\_TIME\_LIMIT\_MS \{3000,8000,15000\}}, \texttt{MIP\_LONG\_WINDOW\_DIVISOR \{4,6,10\}}.
  \item \textbf{CP}: \texttt{CP\_MAX\_TIME\_SECONDS \{3,8,15\}}, \texttt{CP\_LONG\_WINDOW\_DIVISOR \{4,6,10\}}.
\end{itemize}

\end{document}
